%!TEX root = ./main.tex

\section[Money]{Does Any of This Make Money?}

% SECTION TIMING: 8 frames, about 15 minutes. This section closes the loop opened by
% the two money frames in 01_your_phone.tex -- the $50 bill and the flat-revenue /
% rising-spend pair. End on the Starlink question, not on a summary: it is the one
% students will actually go look up afterwards.
%
% VERIFY BEFORE CLASS: every figure in this section moves. Re-check each one and update
% the "as of" line on the slide. A stale number that is visibly dated is acceptable; a
% stale number presented as current is not.

\begin{frame}{The question we opened with}
  \vspace{0.6em}
  \centering
  \begin{tikzpicture}[scale=0.9,transform shape]
    \draw[flow,draw=softgreen,very thick] (0,0.6) -- (5.4,0.6);
    \node[font=\small,anchor=south] at (2.7,0.7) {AT\&T revenue per line: flat};
    \draw[flow,draw=coreorange,very thick] (0,-0.6) -- (5.4,0.35);
    \node[font=\small,anchor=north] at (2.7,-0.7) {what AT\&T spends building it: up};
  \end{tikzpicture}

  \vspace{1.2em}
  \large
  Two lines going different directions.

  \vspace{0.8em}
  \ask{How long can that last, and what happens at the end of it?}
\end{frame}

% Back-of-envelope. BOTH INPUTS ARE AT&T, matching 01_your_phone.tex. An earlier draft
% divided Verizon capex by Verizon connections here while the setup frames used AT&T
% revenue and AT&T spend, so the section changed carriers mid-argument.
%
% SOURCE: AT&T FY2025 ACTUAL consolidated capital expenditures $20.8B (vs $20.3B in
%   FY2024). https://www.prnewswire.com/news-releases/att-reports-strong-fourth-quarter-and-full-year-2025-financial-performance-driven-by-growth-in-converged-fiber-and-5g-customers-302672564.html
% CAPEX, NOT CAPITAL INVESTMENT. AT&T's FY2025 capital investment was $22.0B; the
%   $23-24B figure on the earlier frame is GUIDED capital investment for 2026-2028.
%   Capital investment = capex + cash paid for vendor financing ($1.2B in FY2025).
%   The glossline on this slide exists so students do not read the two as contradictory.
% SOURCE: AT&T wireless subscribers 120.1M at 2025-12-31 (117.9M a year earlier).
%   FY2025 Form 10-K. Breakdown: 91M postpaid, of which 74M phone; 18M prepaid;
%   11M reseller. https://www.sec.gov/Archives/edgar/data/732717/000073271726000120/t-20251231.htm
% Denominator choice: TOTAL wireless subscribers, not postpaid phone. Consolidated
%   capex serves the whole base, so dividing it by the 74M phone lines alone would
%   charge phone customers for capacity built for everyone and give $23/month.
% The division is deliberately crude and the slide says so, in both directions.
\begin{frame}{Roughly where a quarter of your bill goes}
  \vspace{0.3em}
  \centering
  {\large
  \begin{tabular}{@{}r@{\hspace{0.5em}}l@{}}
    \textcolor{coreorange}{\textbf{\$20.8 billion}} & a year of network construction\\[4pt]
    \textcolor{softgray}{$\div$} & \\[4pt]
    \textcolor{ranblue}{\textbf{120 million}} & wireless subscribers\\[6pt]
    \midrule
    $\approx$ \textbf{\$14} & per subscriber per month\\
  \end{tabular}\par}

  \vspace{0.7em}
  \small
  Against a bill of roughly \$50, that is about a quarter of what you pay going
  straight back into towers, spectrum, fiber, and power.

  \vspace{0.4em}
  {\scriptsize\textcolor{softgray}{Crude on purpose, and it cuts both ways:
   consolidated capex also builds fiber and business networks, which pushes the true
   wireless figure down; but 120 million counts prepaid and reseller subscribers
   paying well under \$50, which pushes it up. AT\&T FY2025 capex and year-end 2025
   subscribers.}\par}

  \glossline{The \$23--24 billion from the opening money slide is
    \emph{capital investment} ---
    capex plus vendor financing, and a forecast. This \$20.8 billion is capex actually
    spent in 2025.}
\end{frame}

\begin{frame}{Four national carriers became three}
  \vspace{0.8em}
  \centering
  \begin{tikzpicture}[scale=0.88,transform shape,
      ev/.style={font=\scriptsize,align=center,text width=2.5cm}]
    \draw[flow,draw=softgray,very thick] (-0.6,0) -- (10.4,0);
    \foreach \x/\y in {0/2004, 2.4/2005, 5.6/2009, 9.2/2020} {
      \filldraw[fill=ranblue,draw=ranblue] (\x,0) circle (3pt);
      \node[font=\small\bfseries,anchor=north] at (\x,-0.2) {\y};
    }
    \node[ev,anchor=south] at (0,0.25)   {Cingular buys\\AT\&T Wireless};
    \node[ev,anchor=south] at (2.4,0.25) {Sprint merges\\with Nextel};
    \node[ev,anchor=south] at (5.6,0.25) {Verizon buys\\Alltel};
    \node[ev,anchor=south] at (9.2,0.25) {T-Mobile\\merges Sprint};
  \end{tikzpicture}

  \vspace{1.4em}
  \small
  Years shown are when each deal \textbf{closed}. Brands lingered for years after,
  which is why you still hear some of these names.

  \vspace{0.6em}
  \takeaway{This is a market that consolidates, not one that fragments.}
\end{frame}

\begin{frame}{So are they in trouble, or are they fine?}
  \begin{columns}[T,onlytextwidth]
    \column{0.47\textwidth}
      \large\textbf{\textcolor{softgreen}{Fine}}

      \vspace{0.5em}
      \small
      \begin{itemize}\itemsep5pt
        \item Three players, high barriers
        \item Spectrum licenses are scarce
        \item Everyone needs a phone
        \item Steady, predictable revenue
      \end{itemize}
    \column{0.47\textwidth}
      \large\textbf{\textcolor{coreorange}{In trouble}}

      \vspace{0.5em}
      \small
      \begin{itemize}\itemsep5pt
        \item Revenue per line is not growing
        \item Every generation costs billions
        \item Customers switch on price
        \item Debt from spectrum and mergers
      \end{itemize}
  \end{columns}

  \vspace{0.9em}
  \qa{Which of these is a network architecture problem?}
     {Most of them. Cheaper coverage, shared infrastructure, and software you can
      upgrade without replacing hardware are all architecture answers to a
      financial question. That is a large part of why O-RAN exists.}
\end{frame}

\begin{frame}{The competitor already in your pocket}
  \begin{columns}[T,onlytextwidth]
    \column{0.50\textwidth}
      \large\textbf{Wi-Fi.}

      \vspace{0.6em}
      \small
      Every hour you spend on home or campus Wi-Fi is an hour the carrier
      does not have to build capacity for.

      \vspace{0.6em}
      That is help --- and it is also a ceiling on what they can charge for
      data you were never going to buy from them.
    \column{0.46\textwidth}
      \vspace{0.4em}
      \ask{When did you last care which one you were on?}

      \vspace{0.6em}
      {\small The answer is usually: when something broke.\par}
  \end{columns}

  \vspace{0.7em}
  \takeaway{A competitor you do not notice is still a competitor.}
\end{frame}

% SOURCE: Starlink direct-to-cell status and the June 2026 market reaction.
% Verizon fell about 7%; T-Mobile touched a 52-week low; T-Mobile's CEO publicly
% called the threat exaggerated. VERIFY and update -- this story is still moving.
\begin{frame}{The competitor in orbit}
  \begin{columns}[T,onlytextwidth]
    \column{0.50\textwidth}
      \large\textbf{Starlink direct-to-cell.}

      \vspace{0.6em}
      \small
      \begin{itemize}\itemsep5pt
        \item Ordinary phones, no special hardware
        \item Started as a carrier partnership
        \item Then began looking like a rival
      \end{itemize}
    \column{0.46\textwidth}
      \large\textbf{June 2026.}

      \vspace{0.6em}
      \small
      On the news that SpaceX might go direct to consumers, carrier
      shares fell sharply in a single session.

      \vspace{0.5em}
      T-Mobile's CEO called the threat exaggerated.
  \end{columns}

\end{frame}

\begin{frame}{What a satellite can and cannot take}
  \vspace{0.4em}
  \begin{columns}[T,onlytextwidth]
    \column{0.47\textwidth}
      \large\textbf{\textcolor{softgreen}{Satellites win}}

      \vspace{0.5em}
      \small
      Where there are no towers and never will be: oceans, deserts,
      mountains, disasters.

      \vspace{0.5em}
      Coverage is their whole argument.
    \column{0.47\textwidth}
      \large\textbf{\textcolor{coreorange}{Towers win}}

      \vspace{0.5em}
      \small
      Where the people are. A stadium needs capacity in one small place,
      and capacity comes from cells close together.

      \vspace{0.5em}
      One satellite covers a huge area --- and everyone in it shares.
  \end{columns}

  \vspace{0.9em}
  \takeaway{Coverage and capacity are different products. So far they are
    losing to each other on different ground.}
\end{frame}

\begin{frame}{The question to leave with}
  \vspace{1.0em}
  \centering
  \Large
  Every architecture change we walked through\\
  was somebody trying to serve more people\\
  for less money per person.

  \vspace{1.4em}
  \normalsize
  Splitting boxes. \\Moving voice onto IP. \\Turning the core into software.\\
  Opening the RAN so you can buy the pieces separately.

  \vspace{1.4em}
  \ask{If that pressure never stops, what does the network look like next?}
\end{frame}
